% ============================================================================
% APPENDIX A: APSIM COMPONENTS REFERENCE
% ============================================================================

\chapter{APSIM Components Reference}
\label{app:components}

\section{Global Components}

\subsection{Clock}

\apsimterm{Clock} advances simulation by 1 day. User specifies:
\begin{itemize}
  \item Start date (e.g., 1900-01-01)
  \item End date (e.g., 2000-12-31)
  \item Timestep (always 1 day in APSIM)
\end{itemize}

Role: Synchronizes all other components to daily cycle.

\subsection{Weather}

\apsimterm{Weather} reads external file (\filename{.met} format) with daily:
\begin{itemize}
  \item MaxT, MinT (°C)
  \item Rainfall (mm)
  \item Radiation (MJ/m²/day)
  \item RelativeHumidity (\%)
  \item WindSpeed (m/s)
\end{itemize}

Role: Provides atmospheric drivers; calculates PET (Potential Evapotranspiration).

\subsection{SoilArbitrator}

\apsimterm{SoilArbitrator} coordinates water and nutrient availability among competing processes.

Role: Ensures consistency when multiple components access soil (plant uptake, evaporation, nitrification).

\subsection{MicroClimate}

\apsimterm{MicroClimate} partitions solar radiation and calculates surface energy balance.

Role: Computes interception, albedo, near-surface temperature; feeds into soil heat flux.

\section{Zone Components}

\subsection{Soil}

\apsimterm{Soil} represents multi-layer soil profile (typically 7--10 layers).

\subsubsection{Sub-components}

\begin{description}
  \item[\apsimterm{Physical}] Bulk density, water characteristics (LL, DUL, SAT), saturated conductivity for each layer; includes crop-specific (\apsimterm{SoilCrop}) parameters (LL, KL, XF)
  
  \item[\apsimterm{SoilWater}] (WaterBalance model) Calculates daily water budget (infiltration, drainage, evaporation, plant uptake); tracks water stress factors
  
  \item[\apsimterm{Organic}] Organic matter pools (Fresh Organic Matter, Humic Pool) and C:N ratios; mineralization kinetics
  
  \item[\apsimterm{Chemical}] pH, electrical conductivity, exchangeable sodium percentage; available for future expansion (salinity, boron, etc.)
  
  \item[\apsimterm{Nutrient}] (or \apsimterm{Organic}) Organic N pool; mineralization-immobilization rates
  
  \item[\apsimterm{Solutes}] Individual solute pools: NO\textsubscript{3}, NH\textsubscript{4}, Urea; movement with water; leaching
  
  \item[\apsimterm{Temperature}] Soil temperature by depth; heat diffusion; important for mineralization, nitrification rates
\end{description}

\subsection{Crop (e.g., Wheat)}

\apsimterm{Wheat} (or other crop model) simulates plant processes:

\begin{description}
  \item[\apsimterm{Phenology}] Thermal time accumulation; Zadok stages; development rate parameters by stage
  
  \item[\apsimterm{Leaf}] LAI (Leaf Area Index) progression; specific leaf area (SLA); leaf senescence
  
  \item[\apsimterm{Root}] Root depth progression; water and N extraction coefficients; root biomass
  
  \item[\apsimterm{Biomass}] Daily photosynthesis (RUE: Radiation Use Efficiency); respiration; senescence; partitioning to organs
  
  \item[\apsimterm{Grain}] Grain number determination; grain fill duration; harvest index; grain N content
\end{description}

\subsection{Manager}

\apsimterm{Manager} user-defined control logic (JavaScript or C\#).

Typical rules:
\begin{itemize}
  \item Sowing: IF (date in window AND ESW sufficient AND recent rain) THEN Sow()
  \item Fertilisation: IF (date or stage trigger) THEN Fertiliser.Apply()
  \item Harvest: IF (crop physiologically mature) THEN Harvest()
\end{itemize}

\subsection{Report}

\apsimterm{Report} specifies output variables (acts as filter).

User lists variables to record each day; APSIM records only selected variables (reduces output file size).

\subsection{DataStore}

\apsimterm{DataStore} SQL-like database; receives Report data each day; exports to Excel/CSV at simulation end.

\section{Variable Naming Convention}

APSIM variables use dot notation:

\begin{verbatim}
[ComponentName].SubComponent.Variable

Examples:
[Wheat].LAI                            # Leaf Area Index
[Wheat].Grain.Total.Wt                 # Total grain mass
[Soil].ESW                             # Extractable soil water
[Soil].SoilWater.Evaporation           # Daily evaporation
[Weather].Rainfall                     # Daily rainfall
[Clock].Today                          # Current date
\end{verbatim}

\textbf{Note:} Component names are enclosed in square brackets when used in Report variable lists.

% ============================================================================
% END APPENDIX A
% ============================================================================
